\section{Gauge bridge: Berry/Wilczek--Zee connections as gauge structure}
\label{sec:gauge}

\subsection{Per-wavepacket band isolation and polarization transport}
\label{subsec:per-wavepacket-isolation}

The complex-envelope description is applied per band-isolated excitation, not to a globally
coherent carrier sector spanning macroscopic scales.
Each wavepacket is treated as narrowband around its own local carrier angular frequency $\omega_0$,
determined by the linearized carrier operator about the slowly varying background at the
wavepacket's location.
A band-isolated wavepacket can be written locally as
\begin{equation}
  \vecu(x,t)\;\approx\;\mathrm{Re}\Bigl\{
    a(x,t)\,\bm{\epsilon}(x,t)\,e^{\ii\theta(x,t)}
  \Bigr\},
  \qquad \theta(x,t)=\phi(x,t)-\omega_0 t,
  \label{eq:quasi-monochromatic}
\end{equation}
where $a(x,t)\ge 0$ is the slowly varying envelope amplitude, $\bm{\epsilon}(x,t)\in\C^{N}$ is a
normalized polarization state ($\bm{\epsilon}^\dagger\bm{\epsilon}=1$), and $\omega_0$ is the
wavepacket's local carrier angular frequency.

A central point is that polarization transport is not unique: the same physical state can be
represented in many local phase conventions.
Geometric phase formalizes this redundancy as a connection on a bundle of polarization states over
the spacetime manifold parameterizing the wavepacket's trajectory.

\subsection{Rank-1 Berry connection and electromagnetic identification}

For a single isolated polarization mode, represent the local state by a normalized vector
$|u(x)\rangle\in\C^{N}$ that depends smoothly on spacetime position $x^\mu$
(with $x^\mu=(ct,x^1,x^2,x^3)$ built from the intrinsic coordinates and a chosen branch speed $c$).
Define the Berry connection
\begin{equation}
  a_\mu(x) \;:=\; \ii\,\langle u(x)\,|\,\partial_\mu u(x)\rangle.
  \label{eq:berry-connection}
\end{equation}
Under a local phase choice $|u(x)\rangle\mapsto e^{-\ii\chi(x)}|u(x)\rangle$,
\begin{equation}
  a_\mu \;\mapsto\; a_\mu + \partial_\mu \chi,
  \label{eq:berry-gauge}
\end{equation}
so only gauge-invariant quantities carry physical meaning.
The Berry curvature is
\begin{equation}
  f_{\mu\nu} \;:=\; \partial_\mu a_\nu - \partial_\nu a_\mu,
  \label{eq:berry-curvature}
\end{equation}
and the holonomy around a closed loop $\Gamma$ is
\begin{equation}
  \gamma_\Gamma
  = \oint_\Gamma a_\mu\,\dd x^\mu
  = \int_{\Sigma:\,\partial\Sigma=\Gamma} f_{\mu\nu}\,\dd\Sigma^{\mu\nu}.
  \label{eq:berry-holonomy}
\end{equation}

\paragraph{Electromagnetic reading: four-potential and field tensor.}
In the $U(1)$ trace sector (\Cref{subsec:u3-decomp}) the Berry connection
\eqref{eq:berry-connection} \emph{is} the electromagnetic four-potential, $A_\mu\equiv a_\mu$,
and its curvature \eqref{eq:berry-curvature} \emph{is} the field tensor $F_{\mu\nu}\equiv f_{\mu\nu}$,
with $E_i=F_{0i}$ (time--space) and $B_i=\tfrac12\epsilon_{ijk}F_{jk}$ (space--space).
Two consequences are automatic, requiring no field dynamics:
\begin{enumerate}
  \item the gauge freedom \eqref{eq:berry-gauge} is exactly the unobservable choice of carrier phase
  reference, and
  \item the homogeneous Maxwell equations are the Bianchi identity $\partial_{[\lambda}F_{\mu\nu]}=0$,
  which holds identically because $F=\mathrm{d}A$.
\end{enumerate}
A nonzero field ($F\neq 0$) requires the phase to be non-integrable: the polarization state
$|u\rangle$ must genuinely rotate within $\C^{N}$ ($N\ge 2$), not merely carry an overall phase,
since a pure gradient $A_\mu=\partial_\mu\chi$ is pure gauge.
What this kinematics does \emph{not} fix is field \emph{dynamics}: the inhomogeneous equations
$\mathrm{d}{\star}F=J$ hold only if the substrate furnishes a Maxwell action
$-\tfrac14\int F\wedge{\star}F$, a separate open derivation.

\subsection{Eigenbundle picture and adiabaticity requirement}
\label{subsec:berry-eigenbundle}

The connection \eqref{eq:berry-connection} is not attached to an arbitrary amplitude/phase split.
The physical choice is $|u(x)\rangle$ as a locally defined eigenvector of the linearized carrier
operator about the slowly varying background at $x$.
As $x$ varies, these eigenvectors form an eigenbundle, and the Berry/Wilczek--Zee connection is the
natural connection on that bundle --- the same eigenbundle picture that underlies Berry-phase effects
in crystalline electronic systems \citep{XiaoChangNiu2010Berry}, here applied to the polarization
eigenbundle of the mechanical carrier.

A clean rank-1 ($U(1)$) reduction requires dynamics to remain in a single isolated band.
If $\Delta\omega_{\mathrm{gap}}(x)$ is the local gap to the nearest competing band and
$\Omega_{\mathrm{mix}}(x)$ is the non-adiabatic mode-conversion rate induced by background
gradients, the adiabatic condition
\begin{equation}
  \Omega_{\mathrm{mix}}(x) \;\ll\; \Delta\omega_{\mathrm{gap}}(x)
  \label{eq:adiabatic-condition}
\end{equation}
must hold on the envelope scale.
When \eqref{eq:adiabatic-condition} fails, the correct reduction retains the relevant
multi-dimensional subspace (the Wilczek--Zee description).
The prestress parameter $\alpha$ and wavelength regime jointly control branch separation and
therefore the validity of the Abelian gauge description.

\subsection{Wilczek--Zee connection for the degenerate triplet}
\label{subsec:wz-connection}

For an $n$-dimensional near-degenerate subspace, choose an orthonormal frame
$\mathbf{E}(x)\in\C^{N\times n}$ with columns $|u_a(x)\rangle$ spanning the subspace
($\mathbf{E}^\dagger\mathbf{E}=I_n$).
The Wilczek--Zee connection is
\begin{equation}
  \mathcal{A}_\mu(x) \;:=\; \ii\,\mathbf{E}(x)^\dagger\,\partial_\mu \mathbf{E}(x)
  \;\in\; \mathfrak{u}(n).
  \label{eq:wz-connection}
\end{equation}
Under $\mathbf{E}\mapsto \mathbf{E}U(x)$ with $U(x)\in U(n)$,
\begin{equation}
  \mathcal{A}_\mu \;\mapsto\; U^\dagger \mathcal{A}_\mu U + \ii\,U^\dagger \partial_\mu U,
  \label{eq:wz-gauge}
\end{equation}
and the non-Abelian curvature is
\begin{equation}
  \mathcal{F}_{\mu\nu}
  \;:=\; \partial_\mu \mathcal{A}_\nu - \partial_\nu \mathcal{A}_\mu
    - \ii[\mathcal{A}_\mu, \mathcal{A}_\nu].
  \label{eq:wz-curvature}
\end{equation}
For $n=3$ (the lateral triplet, \Cref{sec:color}), $\mathcal{A}_\mu\in\mathfrak{u}(3)$ decomposes
into a $U(1)$ trace part and a traceless $SU(3)$ part, with the commutator term in
\eqref{eq:wz-curvature} carrying the self-coupling of the non-Abelian sector.

\paragraph{Gauge-covariant envelope dynamics.}
Projecting the full substrate equations of motion onto the retained $n$-dimensional subspace produces
a gauge-covariant derivative acting on the slow envelope $\Psi\in\C^n$:
\begin{equation}
  \mathcal{D}_\mu \Psi \;=\; \bigl(\partial_\mu + \ii\,\mathcal{A}_\mu\bigr)\Psi,
  \label{eq:cov-deriv}
\end{equation}
which transforms covariantly under local frame rotations.
The coarse-grained slow dynamics then admits a gauge-invariant effective action
\begin{equation}
  S_{\text{eff}}[\Psi,\mathcal{A}]
  = \int \dd^4x\;\Bigl(
    \Psi^\dagger \mathcal{K}(\mathcal{D}) \Psi
    - \kappa\,\mathrm{tr}(\mathcal{F}_{\mu\nu}\mathcal{F}^{\mu\nu})
    + \cdots
  \Bigr),
  \label{eq:Seff}
\end{equation}
where $\mathcal{K}(\mathcal{D})$ is determined by the linearized carrier band and $\kappa$ is an
emergent curvature stiffness.

\subsection{Where the gauge structure lives: spacetime, not the Brillouin zone}
\label{subsec:bz-gauge}

The gauge connection does \emph{not} live in the Brillouin zone.
The dynamical matrix $D(\mathbf{k})$ is the Hessian of a real-valued spring energy: real and symmetric
for every $\mathbf{k}$.
Real symmetric matrices admit real eigenvectors, and a real eigenvector bundle over any
simply-connected BZ region is topologically trivial.
Every plaquette holonomy reduces to the identity and every Berry/Wilczek--Zee curvature vanishes
identically.
We have verified this numerically across a full $\alpha$ sweep over the entire $(k_x,k_y)$ plane:
holonomy is identity to machine precision at every plaquette and every $\alpha$.

The physical gauge connection instead lives over spacetime $(x,t)$.
It is the holonomy of the complex carrier envelope $\Psi\in\C^3$ transported around loops in real
space and time, with the $U(1)$ trace read as the electromagnetic sector.
The Brillouin zone and Bloch's theorem enter only to supply the carrier band and to justify the
narrowband isolation of a wavepacket about its carrier $\mathbf{k}_0$; they are not the base space
of the gauge connection.

The reason the base space is $(x,t)$ and not a spatial slice is that the complex structure itself
originates in the timelike direction: the $U(1)$ phase is the rotation of the real carrier field as
one advances along the time axis (\Cref{sec:bell-interpretation} of Paper~I).
A single spacelike slice is a real snapshot and carries no phase; two adjacent time slices are the
minimum data that fix the $U(1)$.

\subsection{Spinorial holonomy and the $4\pi$ aspect}

For an $n=2$ subspace (a two-level polarization bundle), holonomy lives in a double cover of $SO(3)$.
A closed loop in physical space that induces a $2\pi$ rotation of an internal polarization frame can
therefore produce a sign change of the transported state, while a $4\pi$ loop returns it to itself.
This spinorial $4\pi$ aspect is attributed to Berry/Wilczek--Zee holonomy of the internal
polarization bundle.
Its connection to physical spin-$\tfrac12$ phenomenology is part of the matter bridge (Paper~IV).

\subsection{Gauge bridge open problems}
\label{subsec:gauge-open}

\begin{enumerate}
  \item \textbf{Carrier-mode propagation law.}
  The gauge connection structure ($A_\mu$, $F_{\mu\nu}$, covariant derivative) is established by
  the bridge.
  The outstanding task is to derive that free fluctuations of the Berry/Wilczek--Zee connection
  propagate as massless modes with Maxwell/Yang--Mills dispersion: the carrier-excitation spectrum
  in the appropriate limit should reproduce the massless spin-1 photon/gluon propagation law.

  \item \textbf{Inhomogeneous Maxwell equations.}
  Show that the substrate furnishes the Maxwell action $-\tfrac14\int F\wedge{\star}F$
  (with its emergent Hodge star), so that the inhomogeneous equations $\mathrm{d}{\star}F=J$ hold
  with the correct coupling.

  \item \textbf{Adiabaticity in the nonlinear regime.}
  Verify that the adiabatic condition \eqref{eq:adiabatic-condition} holds in the localized-mode
  parameter regime and quantify the non-adiabatic correction to the holonomy.
\end{enumerate}